\documentclass[11pt]{amsart}
\usepackage[margin=1.2in]{geometry}
\usepackage{amsmath,amssymb,amsthm,booktabs}
\usepackage[colorlinks,linkcolor=blue,citecolor=blue,urlcolor=blue]{hyperref}

\newtheorem{theorem}{Theorem}
\newtheorem{lemma}[theorem]{Lemma}
\newtheorem{corollary}[theorem]{Corollary}
\theoremstyle{remark}
\newtheorem{remark}[theorem]{Remark}

\newcommand{\srg}{\operatorname{srg}}

\title[Hexagons in a Conway 99-graph]{If Conway's 99-graph exists,\\ it has at least 209,781 induced hexagons}
\author{Tom Maciejewski}
\address{Independent researcher, New Jersey, USA}
\email{tommaciejewski4@gmail.com}
\subjclass[2020]{05E30, 05C35, 68V15}
\keywords{strongly regular graph, Conway's 99-graph problem, SAT, proof certificate, Lean}

\begin{document}

\begin{abstract}
Conway's 99-graph problem asks whether a strongly regular graph with
parameters $(99,14,1,2)$ exists. Reimbayev showed that such a graph has at
least $209{,}286$ induced hexagons, with equality exactly when no two disjoint
triangles are joined by exactly two edges; by a theorem of Makhnev, equality
cannot hold. We prove that such a graph has at least $209{,}781$ induced
hexagons, $495$ above Reimbayev's bound. The excess is the number of pairs of
disjoint triangles joined by exactly two edges, for which the best bound we
know of was $96$. The proof has two parts. A counting identity, proved by hand
and in Lean, expresses this number through a local statistic $D(x)\le84$ at
each vertex $x$. A SAT computation shows that $D(x)\le74$ at every vertex. Its
refutation is accepted by two independent proof checkers, one of them formally
verified, and the symmetry breaking it uses is proved sound in Lean as an
abstract theorem; the translation of the problem into the instance is tested,
not formally verified. The existence problem remains open.
\end{abstract}

\maketitle

\section{Introduction}

A graph is \emph{strongly regular} with parameters $(v,k,\lambda,\mu)$ if it
has $v$ vertices, is $k$-regular, and any two adjacent vertices have $\lambda$
common neighbours while any two nonadjacent vertices have $\mu$. Whether an
$\srg(99,14,1,2)$ exists is the second of Conway's five \$1{,}000
problems~\cite{Conway}. It is open, and this paper does not decide it.

Let $\Gamma$ be an $\srg(99,14,1,2)$. Since $\lambda=1$, every edge lies in
exactly one triangle. We call the triangles \emph{lines}. Two disjoint lines
are joined by at most three edges (Lemma~\ref{lem:matching}). Following
Reimbayev~\cite{Reimbayev}, let
\begin{align*}
  n_1&=\text{the number of unordered pairs of disjoint lines joined by exactly three edges},\\
  n_3&=\text{the number of unordered pairs of disjoint lines joined by exactly two edges}.
\end{align*}
The subscripts are his names and do not count edges.

\begin{theorem}\label{thm:main}
Every $\srg(99,14,1,2)$ satisfies $n_3\ge495$ and $n_1\le1221$.
\end{theorem}

Reimbayev proved that the number of induced hexagons in an $\srg(n,k,1,2)$ is
$F(n,k)+n_3$, where $F(n,k)=\frac1{12}nk(k-2)(2k^2-21k+53)$, and he
conjectured that $n_3=0$ always~\cite{Reimbayev}. That is false in general: an
$\srg(243,22,1,2)$ with $n_3=19{,}440$ is known~\cite{Djematene}. For
$(n,k)=(99,14)$ the value of $F$ is $209{,}286$.

\begin{corollary}\label{cor:hexagons}
Every $\srg(99,14,1,2)$ has at least $209{,}781$ induced hexagons.
\end{corollary}

\subsection*{What was known}
The identity $3n_1+n_3=4158$ holds in every $\srg(99,14,1,2)$, so $n_3$ is a
multiple of $3$ between $0$ and $4158$. Makhnev~\cite{Makhnev} proved in 1988
that no $\srg(99,14,1,2)$ satisfies a condition which Reimbayev showed to be
equivalent to $n_3=0$, and Keramatipour~\cite[Theorem~3.4.2]{Keramatipour}
gives a second proof that $n_3\ne0$. So $n_3\ge3$ follows from the refereed
literature; that observation is Reimbayev's. A 2026 preprint of
Kova\v{l}~\cite[Corollary~3.3]{Koval} gives $n_1\le1354$, that is, $n_3\ge96$.
Theorem~\ref{thm:main} improves this to $495$, by a route that uses neither
Makhnev's theorem nor his condition.

\subsection*{Method}
Fix a vertex $x$. Its $14$ neighbours form seven adjacent pairs, and the other
$84$ vertices correspond to the nonadjacent pairs of neighbours
(Section~\ref{sec:rooted}). This well-known picture defines $84$
\emph{transition positions} at $x$, each of which is \emph{canonical} or not.
Let $D(x)$ be the number of canonical ones. Two results give
Theorem~\ref{thm:main}:
\begin{itemize}
\item Theorem~\ref{thm:flags}: $n_1=\frac16\sum_xD(x)$ and
$n_3=\frac12\sum_x(84-D(x))$. This is proved by hand in
Section~\ref{sec:counting} and also in Lean.
\item Theorem~\ref{thm:sat}: $D(x)\le74$ for every vertex $x$. This is a
computation, described in Section~\ref{sec:sat}.
\end{itemize}
Section~\ref{sec:trust} says exactly which steps are machine-checked and which
are not.

\section{The rooted picture}\label{sec:rooted}

Let $\Gamma$ be an $\srg(99,14,1,2)$, let $x$ be a vertex and $N(x)$ its
neighbourhood.

Each $a\in N(x)$ has exactly one neighbour in $N(x)$, the common neighbour of
the adjacent vertices $x$ and $a$. So $N(x)$ induces a perfect matching. Write
$a'$ for the partner of $a$; the lines through $x$ are the seven sets
$\{x,a,a'\}$.

Let $y$ be a vertex outside $\{x\}\cup N(x)$. Then $N(y)\cap N(x)$ is a pair
$\{a,b\}$, and $b\ne a'$, since otherwise the adjacent vertices $a,a'$ would
have the two common neighbours $x$ and $y$. Conversely, if $a,b\in N(x)$ are
not partners, they are nonadjacent and have exactly two common neighbours:
$x$, and one vertex outside $N(x)$, which we call $z_{ab}$. (A common
neighbour inside $N(x)$ would be the partner of both.) So the $84$
outside vertices correspond to the $\binom{14}{2}-7=84$ pairs of non-partners,
and $1+14+84=99$. All edges meeting $\{x\}\cup N(x)$ are now known. Only the
$\binom{84}{2}=3486$ adjacencies among the outside vertices are not.

Fix $a\in N(x)$. The neighbours of $a$ are $x$, $a'$ and the twelve vertices
$z_{ab}$ with $b\in N(x)\setminus\{a,a'\}$. As $N(a)$ also induces a perfect
matching, in which $x$ is matched to $a'$, the twelve vertices $z_{ab}$ are
matched among themselves. This defines an involution $M_a$ without fixed
points on $N(x)\setminus\{a,a'\}$:
\[
  M_a(b)=c \iff z_{ab}\sim z_{ac}.
\]
A \emph{transition position} at $x$ is a pair $(a,\{b,b'\})$ consisting of a
neighbour $a$ and a line $\{x,b,b'\}$ not containing $a$. There are
$14\cdot6=84$ of them. The position is \emph{canonical} if $M_a(b)=b'$, that
is, if $z_{ab}\sim z_{ab'}$. Let $D(x)$ be the number of canonical positions
at $x$, so $0\le D(x)\le84$.

If all four positions involving two lines through $x$ are canonical, the nine
vertices they define induce the Paley graph of order $9$, the unique
$\srg(9,4,1,2)$. So $D(x)$ measures how much of this local pattern is present
at $x$.

\section{From the local statistic to the global count}\label{sec:counting}

\begin{lemma}\label{lem:matching}
The edges between two disjoint lines form a matching. In particular there are
at most three of them.
\end{lemma}

\begin{proof}
If a vertex of the first line had two neighbours in the second, it would form
a second triangle on an edge of the second line.
\end{proof}

\begin{theorem}\label{thm:flags}
In every $\srg(99,14,1,2)$,
\[
  n_1=\frac16\sum_{x}D(x),\qquad n_3=\frac12\sum_{x}\bigl(84-D(x)\bigr).
\]
\end{theorem}

\begin{proof}
A \emph{flag} is a triple $(x,a,b)$ with $a,b\in N(x)$ and
$b\notin\{a,a'\}$. There are $14\cdot12=168$ flags at each vertex, two for
each transition position $(a,\{b,b'\})$. Call a flag canonical if its position
is.

Given a flag, let $z=z_{ab}$, let $T'=\{x,b,b'\}$ be the line on the edge
$xb$, and let $T=\{a,z,y\}$ be the line on the edge $az$. The vertex $y$ is
outside $\{x\}\cup N(x)$: it is not $x$ because $x\not\sim z$, and it is not
in $N(x)$ because the only neighbour of $a$ in $N(x)$ is $a'$, and
$a'\not\sim z$. Hence $T$ and $T'$ are disjoint, and $y=z_{ac}$ for some $c$;
by definition $c=M_a(b)$. The edges $ax$ and $zb$ join $T$ to $T'$. By
Lemma~\ref{lem:matching} the only other possible edge is $yb'$, and since
$y=z_{ac}$ is adjacent to exactly $a$ and $c$ in $N(x)$, it is present if and
only if $c=b'$. So a canonical flag gives a pair of disjoint lines joined by
three edges, and a noncanonical flag gives a pair joined by two.

We count the flags that give a fixed pair. Let $\{T_1,T_2\}$ be joined by
three edges. Choose which line is $T'$ (two ways), then $x\in T'$ (three
ways), then $b\in T'\setminus\{x\}$ (two ways). Let $a$ and $z$ be the vertices
of the other line joined to $x$ and to $b$. Then $a\not\sim b$ because the
joining edges form a matching, and $z$ is the common neighbour of $a$ and $b$
other than $x$. So $(x,a,b)$ is a flag giving $\{T_1,T_2\}$, and every such
flag arises in this way: there are $12$. If $\{T_1,T_2\}$ is joined by two
edges, choose which line is $T'$ (two ways); the ends of the two edges in $T'$
must be $x$ and $b$, in either order. There are $4$ such flags.

The canonical flags number $2\sum_xD(x)$ and the others
$2\sum_x(84-D(x))$. Dividing by $12$ and by $4$ gives the two formulas.
\end{proof}

Adding the formulas gives $3n_1+n_3=\frac12\cdot99\cdot84=4158$, which is
Reimbayev's identity~\cite{Reimbayev}. Theorem~\ref{thm:flags} is the
per-vertex form of the double count behind Reimbayev's relation
$3n_1+n_3=\frac14nk(k-2)$: a transition position at $x$ is a line through $x$
together with a further neighbour of $x$.

\begin{proof}[Proof of Theorem~\ref{thm:main}, given Theorem~\ref{thm:sat}]
By Theorem~\ref{thm:sat}, $\sum_xD(x)\le99\cdot74=7326=6\cdot1221$. This sum
equals $6n_1$, so $n_1\le1221$ and $n_3=4158-3n_1\ge495$.
\end{proof}

\section{The computation}\label{sec:sat}

\begin{theorem}\label{thm:sat}
In every $\srg(99,14,1,2)$ and at every vertex $x$, $D(x)\le74$.
\end{theorem}

The proof is a refutation by a SAT solver of the statement ``there is an
$\srg(99,14,1,2)$ with a vertex $x$ such that $D(x)\ge75$''. We describe the
instance and the checks; the encoder that writes it is part of the material
accompanying this paper~\cite{Artifact}.

\subsection*{The instance}
By Section~\ref{sec:rooted}, an $\srg(99,14,1,2)$ with a chosen vertex and a
labelling of its neighbours is the same thing as a choice of the $3486$
adjacencies among the $84$ labelled
outside vertices for which the whole graph on $99$ vertices satisfies the
strongly regular conditions. The instance has one Boolean variable for each of
these adjacencies and clauses expressing those conditions. The $84$ canonical
positions at the root correspond to $84$ of the variables, the adjacencies
$z_{ab}\sim z_{ab'}$, and a sequential-counter encoding~\cite{Sinz} of
``at least $75$ of them are true'' is added.

\subsection*{Symmetry}
The labelling is unique up to permuting the seven lines through $x$ and
swapping the two vertices of each. These relabellings form the group
$G=S_2\wr S_7$ of order $7!\cdot2^7=645{,}120$, which acts on the $3486$
variables and preserves both the strongly regular conditions and the set of
canonical positions. For $91$ elements $g\in G$ (the $7$ swaps within one line,
and, for each of the $21$ pairs of lines, their transposition combined with
each of the $4$ swap patterns on them) the instance contains lex-leader
constraints $X\le_{\mathrm{lex}}gX$~\cite{Crawford}, where $X$ lists the $84$
canonical variables first and then the others, each block in increasing order
of its slots. With these constraints the instance has $1{,}807{,}340$ variables
and $4{,}490{,}236$ clauses.

\subsection*{The refutation}
CaDiCaL~3.0.1~\cite{CaDiCaL} reports the instance unsatisfiable after
about an hour and writes an LRAT proof of $13.6$\,GB. That proof is accepted by
\texttt{cake\_lpr}~\cite{CakeLpr}, a checker whose implementation is itself
formally verified, and by \texttt{lrat-check} from the
\texttt{drat-trim} distribution~\cite{DratTrim}. The proof was streamed to both
checkers as it was written; its SHA-256 digest is recorded, so a rerun can be
compared with it. The weaker instance asserting $D\ge76$ is also
refuted, in the same way as the instance for $D\ge75$. The instances
asserting $D\ge77$ and $D\ge78$ are refuted by Glucose~4.2~\cite{Glucose} with DRAT proofs checked by
\texttt{drat-trim} and, after conversion to LRAT, by \texttt{cake\_lpr}
(Section~\ref{sec:repro}).

\subsection*{From one root to every vertex}
The refutation concerns the fixed labelled root. If some $\srg(99,14,1,2)$ had
a vertex $v$ with $D(v)\ge75$, the labelling of Section~\ref{sec:rooted} with
$x=v$ would be a satisfying assignment, after applying an element of $G$ to
meet the lex-leader constraints. So the refutation gives $D(v)\le74$ for every
vertex of every such graph.

\section{What is verified}\label{sec:trust}

A computer-assisted proof is as strong as its weakest link, so we list the
links.

\subsection*{Proved in Lean}
The following are theorems of a Lean~4 development~\cite{Lean4,Mathlib}, and
each depends only on Lean's three standard axioms.
\begin{itemize}
\item The rooted reduction of Section~\ref{sec:rooted}, in both directions,
including the statement that any vertex of a given graph can be taken as the
root.
\item The action of $G$ on labels and on adjacency variables, and the
soundness of lex-leader constraints: for a group acting on the variables and
an invariant set of constraints, adding $X\le_{\mathrm{lex}}gX$ for any set of
group elements and one fixed order preserves satisfiability.
\item Theorem~\ref{thm:flags}, and the derivation of the bounds on $n_1$ and
$n_3$ from a bound on $D(x)$ at every vertex (for the rung $D\le77$; the
same arithmetic gives the other rungs).
\end{itemize}

\subsection*{Checked by proof checkers}
That the instance, as a file, is unsatisfiable.

\subsection*{Tested, not proved}
That the file says what Section~\ref{sec:sat} claims. The instance is produced
by a Python program, and the following facts about it are supported by tests
and exhaustive finite checks but not by formal proofs.
\begin{itemize}
\item The clauses express the strongly regular conditions and the constraint
$D\ge75$. The $84$ counted variables were compared with the definition of the
canonical positions. The generic part of the encoder gives satisfiable
instances on five parameter sets for which strongly regular graphs are known;
the skeleton, lex-leader and counter layers were tested separately.
\item The emitted lex-leader clauses express $X\le_{\mathrm{lex}}gX$.
\item The permutations used are the ones that Lean proves to form the group
action. They were compared entry by entry with tables exported from Lean, and
the invariance of the canonical positions was checked for all $645{,}120$
group elements. The invariance of the strongly regular conditions under $G$ is
tested, not proved; the Lean theorem takes it as a hypothesis.
\end{itemize}
Closing this gap needs a verified encoder, which we have not written. It is
the standard gap in certificate-based
combinatorics~\cite{Pythagorean}.

\subsection*{Cited}
Reimbayev's hexagon formula, used only for Corollary~\ref{cor:hexagons}.

\section{Remarks}

\begin{remark}[Weaker bounds]
Table~\ref{tab:ladder} shows what weaker hypotheses give by the same
arithmetic. The hypothesis $D(x)\le83$ says that no vertex is canonical in all
$84$ positions; its negation has its own certified refutation. The value
$D(x)=83$ is impossible for an elementary reason: if the position
$(a,\{b,b'\})$ is not canonical then neither is $(a,\{c,c'\})$, where
$c=M_a(b)$. Instances asserting $D\ge81$, $D\ge79$, $D\ge78$, $D\ge77$ and
$D\ge76$ were refuted on the way to $D\ge75$. We did not certify lower
rungs: a run on the instance for $D\ge74$ was stopped after four hours.
\end{remark}

\begin{table}[h]
\centering
\begin{tabular}{lrrr}
\toprule
Hypothesis at every vertex & $n_3\ge$ & $n_1\le$ & hexagons $\ge$\\
\midrule
none (Makhnev and Reimbayev) & $3$ & $1385$ & $209{,}289$\\
none (Kova\v{l}~\cite{Koval}) & $96$ & $1354$ & $209{,}382$\\
$D(x)\le83$ & $51$ & $1369$ & $209{,}337$\\
$D(x)\le82$ & $99$ & $1353$ & $209{,}385$\\
$D(x)\le77$ & $348$ & $1270$ & $209{,}634$\\
$D(x)\le76$ & $396$ & $1254$ & $209{,}682$\\
$D(x)\le75$ & $447$ & $1237$ & $209{,}733$\\
$D(x)\le74$ & $495$ & $1221$ & $209{,}781$\\
\bottomrule
\end{tabular}
\medskip
\caption{Bounds for an $\srg(99,14,1,2)$ from bounds on $D$.}\label{tab:ladder}
\end{table}

\begin{remark}[How far this is from the problem]
The bound $D(x)\le74$ excludes the values $75,\dots,84$ of $D(x)$, and
$n_3\ge495$ excludes about an eighth of the range $[0,4158]$. At every vertex at
least ten of the $84$ positions are non-canonical, so the local pattern
which Reimbayev's conjecture would force everywhere is far from complete. The
computation does not suggest that the remaining cases are within reach of
the same method.
\end{remark}

\begin{remark}[Related work]
The rooted picture of Section~\ref{sec:rooted} is not new. It appears in SAT
attacks on the problem by Keramatipour~\cite{Keramatipour} and
Selub~\cite{Selub}, and again in recent work of Thakkar and
Severini~\cite{ThakkarSeverini}. What is added here is its formal proof, the
symmetry breaking over the stabiliser of the root with a proof of soundness,
and the per-vertex statistic $D$. Reimbayev has classified the induced
subgraphs of order six~\cite{Reimbayev6} and seven~\cite{Reimbayev7} in these
graphs.
Automorphisms of a putative graph are severely restricted~\cite{BehbahaniLam};
we do not use this.
\end{remark}

\section{Reproducibility}\label{sec:repro}

The encoder is written in Python~3 with \texttt{python-sat}~1.9.dev7. The
solvers are CaDiCaL~3.0.1 and Glucose~4.2, and the checkers are
\texttt{cake\_lpr} (commit \texttt{a36874a8}) and \texttt{drat-trim} with
\texttt{lrat-check} (commit \texttt{2e3b2dc0}). The Lean development uses
Lean~4.32.2 with the matching Mathlib release. The accompanying
material~\cite{Artifact} contains the encoder, the Lean sources, the commands
that regenerate every instance and check every proof, and a ledger with the
SHA-256 digests of the instances and proofs. The instance behind
Theorem~\ref{thm:sat} has digest
\begin{center}
\texttt{41a769206eb067b42534c9402a5b01f4af61f4c0dd181242519d59086efcf598},
\end{center}
and its LRAT proof has digest
\begin{center}
\texttt{fdbe2e91b1b8780f1f9e91e2798a6ac68268f676ffc032d833d30973d35c42cb}.
\end{center}
The instances for $D\ge78$ and $D\ge77$, and their DRAT and LRAT proofs, were
regenerated from the material and agree byte for byte with the digests in the
ledger.

\begin{thebibliography}{99}

\bibitem{CaDiCaL}
A.~Biere, T.~Faller, K.~Fazekas, M.~Fleury, N.~Froleyks and F.~Pollitt,
CaDiCaL 2.0, in \emph{CAV 2024}, Lecture Notes in Comput. Sci. 14681,
Springer, 2024, 133--152.

\bibitem{Djematene}
D.~Djematene, Strongly regular graphs: the automorphism group of a Conway
99-graph, new graphs with parameters $(243,22,1,2)$, and a 2-adic feasibility
condition, preprint, Zenodo, 2026,
\href{https://doi.org/10.5281/zenodo.23305123}{doi:10.5281/zenodo.23305123}.

\bibitem{Glucose}
G.~Audemard and L.~Simon, Predicting learnt clauses quality in modern SAT
solvers, in \emph{Proc. IJCAI 2009}, 399--404.

\bibitem{BehbahaniLam}
M.~Behbahani and C.~Lam, Strongly regular graphs with non-trivial
automorphisms, \emph{Discrete Math.} \textbf{311} (2011), 132--144.

\bibitem{Conway}
J.~H. Conway, Five \$1,000 problems (update 2017), On-Line Encyclopedia of
Integer Sequences, entry A248380,
\url{https://oeis.org/A248380/a248380.pdf}.

\bibitem{Crawford}
J.~M. Crawford, M.~L. Ginsberg, E.~M. Luks and A.~Roy, Symmetry-breaking
predicates for search problems, in \emph{Proc. KR'96}, Morgan Kaufmann, 1996,
148--159.

\bibitem{Pythagorean}
M.~J.~H. Heule, O.~Kullmann and V.~W. Marek, Solving and verifying the Boolean
Pythagorean triples problem via cube-and-conquer, in \emph{SAT 2016}, Lecture
Notes in Comput. Sci. 9710, Springer, 2016, 228--245.

\bibitem{Keramatipour}
A.~Keramatipour, Approaching the Conway-99 problem using SAT solvers, M.Phil.
dissertation, University of Cambridge, 2023; arXiv:2604.23037.

\bibitem{Artifact}
T.~Maciejewski, Hexagons in a Conway 99-graph: code, certificates and Lean
sources, 2026, \url{https://github.com/tommac/conway99-hexagons}.

\bibitem{Koval}
A.~Kova\v{l}, On the local structure of a hypothetical strongly regular graph
with parameters $(99,14,1,2)$, preprint, Zenodo, 2026,
\href{https://doi.org/10.5281/zenodo.23281436}{doi:10.5281/zenodo.23281436}.

\bibitem{Makhnev}
A.~A. Makhnev, Strongly regular graphs with $\lambda=1$, \emph{Math. Notes}
\textbf{44} (1988), 847--850.

\bibitem{Mathlib}
The mathlib Community, The Lean mathematical library, in \emph{Proc. CPP
2020}, 367--381.

\bibitem{Lean4}
L.~de Moura and S.~Ullrich, The Lean~4 theorem prover and programming
language, in \emph{CADE 28}, Lecture Notes in Comput. Sci. 12699, Springer,
2021, 625--635.

\bibitem{Reimbayev}
R.~Reimbayev, The lower bound for number of hexagons in strongly regular
graphs with parameters $\lambda=1$ and $\mu=2$, \emph{e-J. Anal. Appl. Math.}
\textbf{2024} (2024), 1--14; arXiv:2409.10620.

\bibitem{Reimbayev6}
R.~Reimbayev, The subgraphs of order six of the family of strongly regular
graphs with parameters $\lambda=1$ and $\mu=2$, preprint, 2025,
arXiv:2508.03377.

\bibitem{Reimbayev7}
R.~Reimbayev, Induced subgraphs of order seven and their frequencies in
$\srg(n,k,1,2)$, preprint, 2026, arXiv:2608.19410.

\bibitem{Selub}
N.~Selub, Conway's 99-graph problem: a Boolean satisfiability approach,
University of Chicago REU paper, 2023,
\url{https://math.uchicago.edu/~may/REU2023/REUPapers/Selub.pdf}.

\bibitem{Sinz}
C.~Sinz, Towards an optimal CNF encoding of Boolean cardinality constraints,
in \emph{CP 2005}, Lecture Notes in Comput. Sci. 3709, Springer, 2005,
827--831.

\bibitem{CakeLpr}
Y.~K. Tan, M.~J.~H. Heule and M.~O. Myreen, \texttt{cake\_lpr}: verified
propagation redundancy checking in CakeML, in \emph{TACAS 2021}, Lecture Notes
in Comput. Sci. 12652, Springer, 2021, 223--241.

\bibitem{ThakkarSeverini}
A.~Thakkar and S.~Severini, A forced-structure reduction and verifiable bounds
for Conway's 99-graph, preprint, 2026, arXiv:2608.11211.

\bibitem{DratTrim}
N.~Wetzler, M.~J.~H. Heule and W.~A. Hunt, Jr., DRAT-trim: efficient checking
and trimming using expressive clausal proofs, in \emph{SAT 2014}, Lecture
Notes in Comput. Sci. 8561, Springer, 2014, 422--429.

\end{thebibliography}

\end{document}
